\chapter{Spécialiser un modèle sur son domaine}

Les chapitres précédents ont montré les outils. Celui-ci les met bout à bout sur
un cas complet, du corpus vide au modèle qui répond sur votre sujet.

\section{Généraliste ou spécialisé : ce qui change vraiment}

Rien, dans les commandes. Tout, dans le corpus.

\begin{center}
\begin{tabular}{p{4.0cm}p{9.6cm}}
\toprule
\textbf{Corpus} & \textbf{Modèle obtenu} \\
\midrule
Conversations variées & Généraliste : répond correctement à tout, sans exceller
                        nulle part. \\
Documentation d'un logiciel & Répond sur \emph{ce} logiciel : ses commandes, ses
                        messages d'erreur, ses conventions. \\
Textes juridiques d'un pays & Emploie le vocabulaire et les tournures de ce
                        droit-là. \\
Archives d'un service client & Répond dans le ton de la maison, sur ses
                        produits. \\
\bottomrule
\end{tabular}
\end{center}

\begin{nkretenir}
Vous ne réglez pas la spécialisation, vous la \emph{nourrissez}. Le temps que
vous passerez à choisir vos textes comptera davantage que tous les réglages
réunis.
\end{nkretenir}

\section{Un cas complet, de bout en bout}

Prenons un exemple concret : spécialiser un modèle sur la documentation de votre
propre logiciel, pour qu'il réponde aux questions de vos utilisateurs.

\subsection{Étape 1 — écrire les questions de référence}

\textbf{Avant tout}, écrivez les questions qui serviront à juger. Dix suffisent.

\begin{nkcode}[]{questions.txt}
Comment reprendre un entraînement interrompu ?
Que se passe-t-il si un checkpoint est corrompu ?
Quelle est la différence entre un checkpoint et un fichier d'adaptateurs ?
Pourquoi utiliser Vulkan plutôt que DirectX ?
Combien de temps prend une époque complète ?
\end{nkcode}

\begin{nkpiege}
Écrire ces questions \emph{après} avoir vu les réponses du modèle est le plus sûr
moyen de se convaincre d'un progrès qui n'existe pas. On choisit alors, sans le
vouloir, les questions où le modèle brille.
\end{nkpiege}

\subsection{Étape 2 — mesurer l'avant}

\begin{nkcode}[]{La référence}
NKQwen2Ask --gguf=<modele> --question="Comment reprendre un entraînement
interrompu ?" --tokens=64
\end{nkcode}

Sans corpus spécialisé, le modèle répondra quelque chose de générique et de
plausible — souvent faux dans les détails. C'est exactement ce qu'on veut
constater, et conserver.

\subsection{Étape 3 — construire le corpus}

Reprenez la méthode du chapitre~3 : découper, interroger, répondre.

\begin{nkcode}[]{corpus-doc.txt — un extrait}
Question: Comment reprendre un entraînement interrompu ?
Reponse: Relancez exactement la même commande. Le programme cherche les
checkpoints, retient le plus avancé qui soit intact, et repart du pas et
de l'exemple où il s'était arrêté.

Question: Que se passe-t-il si un checkpoint est corrompu ?
Reponse: Son empreinte ne correspond plus, il est refusé et ignoré. La
reprise se fait depuis le second checkpoint, sans perdre de pas.

Question: Quelle est la différence entre un checkpoint et un fichier
d'adaptateurs ?
Reponse: Le checkpoint pèse 231 Mo et sert à reprendre un entraînement :
il contient les moments d'Adam et la position dans le corpus. Le fichier
d'adaptateurs pèse 77 Mo, ne contient que ce qui a été appris, et sert à
utiliser ou partager le modèle.
\end{nkcode}

\textbf{Combien en faut-il ?} Pour une documentation, comptez trois à cinq
questions par page. Une documentation de 200 pages donne donc de 600 à
1\,000 paires — de quoi occuper la machine une journée.

\subsection{Étape 4 — vérifier le corpus}

\begin{nkcode}[]{Trente secondes bien investies}
NKQwen2Train --gguf=<modele> --corpus=corpus-doc.txt --steps=2 --fresh
\end{nkcode}

Lisez la ligne \texttt{corpus : N paires}. Si $N$ ne correspond pas à ce que
vous avez écrit, corrigez maintenant, pas dans dix heures.

\subsection{Étape 5 — entraîner}

\begin{nkcode}[]{L'entraînement réel}
NKQwen2Train --gguf=<modele> \
             --corpus=corpus-doc.txt \
             --out=modele-doc \
             --lr=1e-4 \
             --save-every=100 \
             --valid-every=200
\end{nkcode}

Laissez tourner. Surveillez deux choses, et deux seulement :

\begin{itemize}
  \item la \textbf{perte moyenne} doit descendre ;
  \item la \textbf{perte de contrôle} doit descendre aussi.
\end{itemize}

\subsection{Étape 6 — mesurer l'après}

Mêmes questions, mêmes réglages, adaptateurs en plus :

\begin{nkcode}[]{Le verdict}
NKQwen2Ask --gguf=<modele> --adapters=modele-doc.nkla \
           --question="Comment reprendre un entraînement interrompu ?" \
           --tokens=64
\end{nkcode}

\begin{nkretenir}
Le seul jugement qui vaille est celui-ci : sur \emph{vos} questions, préparées
\emph{avant}, la réponse est-elle devenue plus juste ? Pas plus longue, pas plus
jolie : plus juste.
\end{nkretenir}

\section{Quand ça ne marche pas}

\subsection{La perte ne descend pas}

\begin{center}
\begin{tabular}{p{4.6cm}p{9.0cm}}
\toprule
\textbf{Cause probable} & \textbf{Vérification} \\
\midrule
Corpus mal formé & La ligne \texttt{corpus : N paires} donne un $N$ bien
                   inférieur au nombre attendu. \\
Taux d'apprentissage trop faible & Essayez \texttt{-{}-lr=3e-4}. \\
Taux trop fort & La perte oscille violemment ou affiche \texttt{nan}. Divisez
                   par dix, et relancez avec \texttt{-{}-fresh}. \\
Exemples trop longs & Ils sont sautés en silence. Augmentez
                   \texttt{-{}-max-seq}, ou raccourcissez vos réponses. \\
\bottomrule
\end{tabular}
\end{center}

\subsection{La perte descend mais les réponses n'ont pas changé}

C'est le cas le plus fréquent, et le plus mal interprété. Deux explications
possibles :

\begin{enumerate}
  \item \textbf{Vous n'avez pas fait assez de pas.} Quelques centaines déplacent
        le style ; il en faut des milliers pour le fond.
  \item \textbf{Le modèle savait déjà.} Si vos questions portent sur des
        généralités, il y répondait correctement avant. Interrogez-le sur ce
        que \emph{seule} votre documentation contient — un nom de commande, un
        message d'erreur, une convention interne.
\end{enumerate}

\subsection{Les réponses sont devenues moins bonnes}

C'est le \textbf{surapprentissage} : le modèle a tellement bien mémorisé votre
corpus qu'il a désappris à répondre au reste. Signe caractéristique : la perte
d'apprentissage continue de descendre pendant que celle de contrôle remonte.

\begin{nkpiege}
Le remède n'est pas de « moins entraîner » au hasard. Prenez le checkpoint
\emph{d'avant} la remontée — c'est précisément pour cela qu'on sauvegarde
régulièrement. Copiez de temps en temps un checkpoint ailleurs, avec le numéro
du pas dans son nom : vous pourrez revenir en arrière.
\end{nkpiege}

\section{Aller plus loin : plusieurs domaines}

Rien ne vous empêche d'entraîner plusieurs jeux d'adaptateurs sur le même modèle
de base :

\begin{center}
\begin{tabular}{ll}
\toprule
\texttt{modele-doc.nkla} & répond sur votre documentation \\
\texttt{modele-juridique.nkla} & répond sur vos contrats \\
\texttt{modele-support.nkla} & répond dans le ton du service client \\
\bottomrule
\end{tabular}
\end{center}

Chacun pèse 77~Mo, et le modèle de base de 4,4~Go reste partagé. Vous changez de
spécialité en changeant un argument :

\begin{nkcode}[]{Un modèle, trois spécialités}
NKQwen2Ask --gguf=<modele> --adapters=modele-doc.nkla       --question="..."
NKQwen2Ask --gguf=<modele> --adapters=modele-juridique.nkla --question="..."
\end{nkcode}

\begin{nkretenir}
C'est là tout l'intérêt du LoRA : le gros fichier ne bouge jamais, seules les
petites pièces changent. Distribuer une spécialité coûte 77~Mo, pas 4,4~Go.
\end{nkretenir}

\begin{nkexo}
Choisissez un domaine que vous connaissez et dont le modèle ne peut rien savoir
— le règlement intérieur de votre école, la documentation d'un projet à vous.
Écrivez cent paires. Entraînez mille pas. Comparez avant et après sur cinq
questions préparées d'avance. C'est le projet qui transforme ce cours en
compétence.
\end{nkexo}
